\section{Fibrations and the six-dimensional interpretation}
\label{sec:fibrations}

\subsection{The K3-fiber vector}
\label{subsec:k3-fibration}

The two vectors imposing the joint constraint have a fibration
interpretation. The functional $-e_3^*$ on the fan takes value
$1$ on $v_5$, $-1$ on $v_3,v_7,v_8$ and zero elsewhere;
no resolved cone contains rays of both signs. It defines a K3
fibration
\begin{equation}
 \rho:\wideX_9\longrightarrow\mathbb P^1,\qquad
 [\mathrm{K3}]=D_5\sim D_3+D_7+D_8.
 \label{eq:k3-fibration}
\end{equation}
The generic fiber has Picard rank three. Its classes
$(D_2|,D_6|,E|)$ have Gram matrix
\begin{equation}
 \begin{pmatrix}-2&1&0\\1&0&1\\0&1&-2\end{pmatrix},\qquad
 \operatorname{Pic}(\mathrm{K3})\simeq U\oplus\langle-4\rangle.
 \label{eq:k3-lattice}
\end{equation}
The middle class is an elliptic fiber and the other two are
disjoint sections. The horizontal surface $D_6$ is a K3 blown
up once; its exceptional curve is $C_1$. In contrast, $C_2$
is a section of $\rho$.

In the K3-fibration heterotic/type-II dictionary, $D_5$ is the
heterotic dilaton vector multiplet
\cite{Kachru:1995wm,Klemm:1995tj,Aspinwall:1995vk,Antoniadis:1995vz}.
Since $D_5^2=0$, its norm is
\begin{equation}
 \mathcal N_{D_5}=\frac{\operatorname{vol}(\mathrm{K3})^2}{\mathcal V}.
 \label{eq:dilaton-norm}
\end{equation}
Thus one of the shared gaugings is a dilaton gauging; the other,
$-D_6$, restricts to the elliptic fiber class in the K3 lattice.
Vertical versus horizontal M2 curves give the usual perturbative
versus nonperturbative interpretation of $C_1$ and $C_2$ in this
dictionary. This interprets the compact charges but does not
independently construct their heterotic states or condensates.

\subsection{The flat elliptic model and its anomalies}
\label{subsec:sixd}\label{subsec:genus-one}

The nef divisor $w=D_4+D_5=D_3+D_6$ has $w^3=0$ and defines
a genus-one fibration over $\mathbb P^2$. Its fiber $F=w^2$
pairs with $(D_0,D_1,D_2,E)$ as $(3,2,1,1)$ and trivially
with the other divisors. There are two rational sections $D_2,E$
and two non-flat points $q_3,q_4$ on a line $L_0$.
Flopping $C_1$ gives a flat elliptic threefold $X_9^{\rm fl}$
over $B'=\operatorname{Bl}_{q_3,q_4}\mathbb P^2$: the fan
projection to the last two coordinates then maps every cone
to a cone. Write $H,E_3,E_4$ for the base hyperplane and
exceptional classes, and $L=H-E_3-E_4$.

The fiber monomials give a Morrison--Park two-section model
\cite{Morrison:2012ei}, specialized so that $L$ carries an
$I_2$ fiber and the sections are disjoint. The latter condition
removes charge-two singlet loci in the normalization where the
ordinary singlet has charge one. The flat six-dimensional
tensor branch has $T=2$, algebra
$\mathfrak{su}(2)\oplus\mathfrak u(1)$, and gauge curve $L$.
The two exceptional base curves support the geometric E-string
sectors \cite{Ganor:1996pc}. Its vector count is
$T+\operatorname{rank}G+1=5$, as required by
\eqref{eq:x9-spectrum-table}.

The charged matter and anomaly coefficient are
\begin{equation}
 10\,\mathbf2_{1/2}+76\,\mathbf1_1,\qquad
 b=\frac{11}{2}H-\frac32E_3-\frac32E_4.
 \label{eq:x9-height-class}
\end{equation}
Here the gravitational hypermultiplet count weights a doublet
by dimension two. Thus $H_{
m tot}=123+20+76=219$,
$V=4$, and $H_{
m tot}-V+29T=273$.
The nonabelian anomaly conditions use $K_{B'}L=L^2=-1$
and give ten fundamentals. In the conventional Abelian
normalization \cite{Kumar:2010ru,Park:2011ji}, the remaining
checks read
\begin{equation}
\begin{aligned}
 -6K_{B'}b&=81=20(1/2)^2+76,\\
 3b^2&=\frac{309}{4}=20(1/2)^4+76,\\
 bL&=\frac52=10(1/2)^2.
\end{aligned}
\label{eq:x9-abelian-anomaly-checks}
\end{equation}
These agree with the geometric count of $76$ isolated $I_2$
fibers away from $L_0$.

The full height, including its exceptional-base components,
is independently fixed by the sections. Put
$P=\pi^*c_1(B')=3w-D_3-D_4$ and use $D_2$ as zero section.
The Shioda and shifted Kaluza--Klein divisors are
\begin{equation}
 \sigma=E-D_2-P+\tfrac12D_8,\qquad Z=D_2+\tfrac12P,
 \qquad -\pi_*(\sigma^2)=b.
 \label{eq:x9-shioda-kk}
\end{equation}
They obey $\pi_*(\sigma Z)=\pi_*(\sigma D_8)=\pi_*(Z^2)=0$
\cite{Grimm:2015gka}. The half-Cartan term supplies the
geometric center identification: in integral charge normalization
the doublets have charge one and the singlets charge two,
giving $U(2)=(SU(2)\times U(1))/\ZZ_2$. These statements
come from the F-theory geometry, not a heterotic bundle tuning.

\subsection{The two E-string roots and the elliptic-circle limit}
\label{subsec:estring-roots}\label{subsec:finite-circle}

On $X_9^{\rm fl}$ the surfaces $S_i=D_i$, $i=3,4$, are rational
elliptic surfaces. Let $o_i=D_2|_{S_i}$, $s_i=E|_{S_i}$,
$F_i=-K_{S_i}$, and define
\begin{equation}
 r_i=s_i-o_i-F_i.
 \label{eq:x9-estring-roots}
\end{equation}
The two sections are disjoint, have square $-1$ and meet the
fiber once. Hence $r_i^2=-2$ and $r_i o_i=r_i F_i=0$:
each is a root in the $E_8(-1)$ orthogonal complement of
$\langle o_i,F_i\rangle$ \cite{Schutt:2009elliptic}.
Their compact images coincide. Indeed $o_3=-C_1$, $o_4=C_2$,
$F_3=F_4=F$ and $s_4-s_3=R$, so
\begin{equation}
 r_4-r_3=R-C_1-C_2=0.
\end{equation}
The integral curve pairings verifying these identities are
printed in Appendix~\ref{app:reproducibility}. This identifies
two local root images, not the entire microscopic heterotic
$E_8$ gauge algebras or a BPS state for every signed root class.

The transition face also determines whether a six-dimensional
limit exists. Write $J=aD_0+bD_1$ with $a,b\geq0$ and
$a+b>0$, and $\tau=JF=3a+2b$. The original and flat chambers
share this face because $JC_1=0$. The cubic gives
\begin{equation}
 \frac7{24}\leq\frac{\mathcal V}{\tau^3}\leq\frac{83}{162}.
 \label{eq:x9-finite-circle-bound}
\end{equation}
To prove the bound, put $x=a/\tau\in[0,1/3]$. Then
\begin{equation}
 \frac{\mathcal V}{\tau^3}
 =\frac7{24}+\frac34x-\frac18x^2-\frac5{12}x^3.
\end{equation}
Its derivative is decreasing and has positive minimum $19/36$;
the endpoint values give \eqref{eq:x9-finite-circle-bound}.
The F-theory limit requires $\tau/\mathcal V^{1/3}\to0$
\cite{Bonetti:2011mw}, so it is excluded on this face.
This is the elliptic 6d-to-5d circle, not the auxiliary 5d-to-4d
circle of Section~\ref{subsec:x9-circle-limit}. Nor does the
bound force the total Planck volume to vanish at fixed fiber
volume: the section $E$ has zero volume on the face, but
$\operatorname{vol}(D_2)=2a^2$.

For completeness the finite-circle charge relation is visible
without taking either limit. Expand
$J=\tau Z+\pi^*j+z\sigma+\phi D_8$; on the face,
$z=2a+b$ and $\phi=-b/2$. The common root image has
$(Zr,\sigma r,D_8r)=(0,-3/2,1)$, so
\begin{equation}
 J(r+F)=\tau-\tfrac32z+\phi=0.
 \label{eq:x9-affine-charge-relation}
\end{equation}
The section-difference class has one unit of this elliptic
Kaluza--Klein charge. A different zero section changes the
divisor coordinates, not the volume bound. At finite radius,
$j$ is not directly the six-dimensional base K\"ahler class:
the tensor dictionary includes Green--Schwarz shifts
\cite{Bonetti:2011mw,Grimm:2015gka}.

\subsection{The status of the heterotic construction}
\label{subsec:heterotic-k3}

Regard $B'$ as $\mathbb F_1$ blown up once on the positive
section. Its heterotic parent has instanton numbers $(11,13)$;
the extra blowup replaces one instanton by a five-brane
\cite{Morrison:1996na,Morrison:1996xf}. The residual data are
\begin{equation}
 (k_1,k_2;N_5)=(11,12;1),\qquad k_1+k_2+N_5=24.
 \label{eq:x9-heterotic-instantons}
\end{equation}
In this presentation the common root above identifies a root
of the E-string over the negative section with one of the
five-brane's E-string. It does not derive their condensate
dynamics.

An integral subdivision of the Newton polytope along
$m_4+2m_1+m_2=0$ has a reflexive common slice with $30$
lattice points. It produces a candidate elliptic heterotic K3
with a distinguished $I_2$ fiber; an exact member has fibers
$I_2+22I_1$. Its equation and coefficient data are given in
Appendix~\ref{app:reproducibility}. The semistable total family,
bundle or spectral-sheaf data and their gluing remain to be
constructed. In particular the massless $U(1)$ requires a
St\"uckelberg check, not merely a bundle commutant
\cite{Cvetic:2015uwu,Honecker:2006qz}. The F-theory spectrum
and root relation therefore constrain a heterotic construction;
they do not complete it.
